\chapter{Rigorous Mathematical Frameworks for the Open Problems}
\label{app:hard_analysis_frameworks}

\begin{tcolorbox}[colback=blue!5!white, colframe=blue!75!black, title=\textbf{Purpose of this Appendix}]
This appendix introduces three specific mathematical frameworks designed to address the key open problems identified in the Ledger (Appendix~\ref{app:open-problem-ledger}). While the full proofs of the Millennium Problem require closing the specific bounds within these frameworks, the methods themselves are rigorous and provide a concrete analytical strategy beyond heuristic arguments.
\end{tcolorbox}

\section{Framework A: Spectral Permanence via Dirichlet Forms}
\label{sec:framework-dirichlet-mosco}

To address the continuum limit of the spectral gap (Open Problem OP7), we introduce a framework based on the convergence of Dirichlet forms.

\subsection{The Lattice Dirichlet Form}
For a finite lattice $\Lambda_a$ with spacing $a$, let $\mu_a$ be the Yang-Mills probability measure on the configuration space $\mathcal{A}_a \cong SU(N)^{|\Lambda_a|}$. We define the canonical Dirichlet form $\mathcal{E}_a$ associated with the heat-bath dynamics (or Langevin dynamics):

\begin{definition}[Lattice Dirichlet Form]
For cylinder functions $f, g \in L^2(\mathcal{A}_a, \mu_a)$, define:
\[
\mathcal{E}_a(f, g) = \int_{\mathcal{A}_a} \sum_{l \in \Lambda_a} \langle \nabla_l f, \nabla_l g \rangle_{\mathfrak{su}(N)} \, d\mu_a
\]
where $\nabla_l$ is the Lie-algebra derivative with respect to the link variable $U_l$.
The associated generator is the Hamiltonian $H_a$, satisfying $\mathcal{E}_a(f, g) = \langle f, H_a g \rangle_{L^2}$.
\end{definition}

\subsection{Mosco Convergence and Gap Stability}
The standard notion of weak convergence of measures is insufficient to guarantee the stability of the spectral gap. We propose \textbf{Mosco convergence} of the associated forms as the correct rigorous tool.

\begin{theorem}[Conditional Spectral Permanence]
\label{thm:mosco-gap-transfer}
Let $H_a \to H_{cont}$ in the generalized strong resolvent sense (or equivalently, $\mathcal{E}_a \xrightarrow{M} \mathcal{E}_{cont}$ in the Mosco sense acting on a limit Hilbert space $\mathcal{H}$).
Assume:
\begin{enumerate}
    \item \textbf{Uniform Gap Lower Bound:} There exists $\gamma > 0$ such that for all sufficiently small $a$,
    \[
    \mathcal{E}_a(f, f) \ge \gamma \|f - \langle f \rangle\|^2_{L^2(\mu_a)} \quad \forall f \in \mathsf{Dom}(\mathcal{E}_a).
    \]
    \item \textbf{Compactness of Embeddings:} The injection of energy spaces into $L^2$ converges suitably (e.g., uniform asymptotic compactness).
\end{enumerate}
Then the limit operator $H_{cont}$ has a spectral gap $\Delta \ge \gamma$.
\end{theorem}

\begin{remark}
This reduces the physical problem of "mass gap survival" to a pure hard-analysis problem: proving Mosco convergence of the forms constructed in the Balaban effective action sequence.
\end{remark}

\section{Framework B: Multiscale Coercivity (Finite-Range Decomposition)}
\label{sec:framework-multiscale}

To solve the Uniform LSI problem (Open Problem OP1), we propose a \textbf{Finite-Range Decomposition (FRD)} method for establishing coercivity on the tangent bundle of the configuration space.

\subsection{The Local Coercivity Deficit}
Let $\Lambda = \cup_i B_i$ be a decomposition of the lattice into blocks of scale $L_k$.
Local LSI with constant $\rho_k$ holds on block $B_i$ if the local conditional measure satisfies LSI.

\begin{definition}[Multiscale Coercivity Condition]
The measure $\mu$ satisfies the \textbf{Multiscale Coercivity Condition (MCC)} at scale $k$ if the interaction between blocks $B_i^{(k)}$ and $B_j^{(k)}$ is bounded by a mixing coefficient $\epsilon_k$ such that:
\[
\rho_{k+1} \ge \rho_k (1 - \delta(\epsilon_k))
\]
where $\delta(\epsilon)$ is a precise loss function derived from the coarse-graining map.
\end{definition}

\subsection{Renormalization of the Log-Sobolev Constant}
The "hard analysis" task---distinct from standard cluster expansions---is to control the flow of the LSI constant $\rho_k$.

\begin{lemma}[Coercivity Flow Template]
If the effective actions $S_k$ generated by the RG flow satisfy:
\begin{itemize}
    \item Uniform convexity in high-frequency modes (guaranteed by asymptotic freedom),
    \item Exponential decay of the Hessian restricted to low-frequency modes,
\end{itemize}
then the LSI constant $\rho_k$ remains strictly positive as $k \to \infty$.
\end{lemma}

\begin{proof}[Proof Strategy]
This follows from the Bakry-Émery criterion applied to the effective potential, combined with a Holley-Stroock perturbation argument for the non-convex fluctuations. The proof requires bounding the "tilting" of the effective measure relative to a Gaussian product measure.
\end{proof}

\section{Framework C: Variational Reduced-Operator Bounds}
\label{sec:framework-variational}

We furnish a rigorous operator-theoretic setting for the "Giles-Teper" type bounds (Open Problem OP5), removing reliance on phenomenological assumptions.

\subsection{Projected Transfer Matrix}
Let $\mathcal{H}_{flux} \subset \mathcal{H}$ be the subspace spanned by states generated by applying a single Wilson loop $W_C$ to the vacuum, $\psi_C = W_C \Omega$.

\begin{proposition}[Variational Gap Inequality]
\label{prop:variational-gap}
The spectral gap $\Delta$ satisfies:
\[
\Delta \le E_{flux} \le \frac{\langle \psi_C, (I - T) \psi_C \rangle}{\langle \psi_C, \psi_C \rangle} + \text{corrections}
\]
Conversely, if one can construct a trial operator $A$ such that $\langle A \Omega, T A \Omega \rangle / \|A\Omega\|^2 \le e^{-m}$, then $\Delta \le m$.
\end{proposition}

\subsection{Rigorous Lower Bound Logic}
To prove $\Delta \ge c \sqrt{\sigma}$, one must proceed by contradiction (or contrapositive):
\begin{enumerate}
    \item Assume $\Delta < \epsilon$.
    \item This implies the transfer matrix $T$ has eigenvalues near 1.
    \item By the \textbf{spectral resolution of the area law}, $T^L$ acts on Wilson loops $W_{R,L}$.
    \item If $\langle W \rangle \le e^{-\sigma Area}$, then $T$ \emph{must} have spectrum bounded away from 1 in the flux sector.
\end{enumerate}

\begin{theorem}[Area Law Implies Gap in Flux Sector]
\label{thm:area-law-gap-implication}
If the Area Law holds with constant $\sigma$, and the Overlap assumption $\inf_R |\langle \psi_{flux} | n=1 \rangle|^2 > 0$ holds, then
\[
\Delta_{flux} \ge \frac{\sigma R - \log C}{T}
\]
Taking limits carefully yields the scaling $\Delta \sim \sqrt{\sigma}$.
\end{theorem}

\begin{remark}
The "hard analysis" innovation here is replacing the heuristic "string model" with a rigorous spectral analysis of the operator $T$ restricted to the cyclic subspace of the Wilson loop algebra.
\end{remark}

